\documentclass[10pt,a4paper]{amsart}
\usepackage[margin=19mm]{geometry}
\usepackage{amsmath,amssymb,mathtools}
\usepackage[scr=rsfs,cal=euler]{mathalfa}
\usepackage{xfrac}
\usepackage[unicode,hidelinks,bookmarks=false]{hyperref}
\newcommand{\ie}{i.e.\ }
\newcommand{\cf}{cf.\ }
\newcommand{\resp}{respectively\ }
\newcommand{\Subsection}{\subsection}
\providecommand{\nobreakdash}{\nobreak\hbox{-}}
\setlength{\emergencystretch}{2em}
\multlinegap0pt
\usepackage{bm}
\usepackage{bbm}
\usepackage{dsfont}
\usepackage{xspace}
\usepackage{cancel}
\usepackage{mathtools}
\usepackage{amsthm}
\makeatletter
\@ifundefined{theorem}{\newtheorem{theorem}{Theorem}}{}
\makeatother
\title[A Novel Algorithm for Representing Positive Semi-Definite Po]{A Novel Algorithm for Representing Positive Semi-Definite Polynomials as Sums of Squares with Rational Coefficients}
\author{Zhenbing Zeng, Yong Huang, Lu Yang, Yongsheng Rao}
\date{Source: arXiv:2510.01568v2 (CC BY 4.0)}
\begin{document}
\maketitle

\noindent\emph{Source and licence.} A Novel Algorithm for Representing Positive Semi-Definite Polynomials as Sums of Squares with Rational Coefficients. Version arXiv:2510.01568v2 (CC BY 4.0), distributed under the Creative Commons Attribution 4.0 International licence, as recorded in the arXiv licence metadata for this exact version and shown on the abstract page https://arxiv.org/abs/2510.01568v2. Full rights notice: licenses/arxiv-2510.01568-CC-BY-4.0.txt.\par\smallskip

\noindent\textit{Editorial scope.} Counted unit: the Theorem whose source label is \texttt{thm-2} in the article (source0.tex lines 293--305, source label \texttt{thm-2}), delivered together with the article's own complete original proof of it (source0.tex lines 307--417, one \texttt{proof} environment of 111 lines). No mathematics is written, repaired or re-derived: statement and proof are the article's own text line for line. Required context retained uncounted: none. Every definition, display and label of those lines is retained unchanged. Omitted from this package and not redistributed: 1--292, 306--306, 418--2110. Nothing inside a retained excerpt is omitted or altered: every retained body is delivered exactly as the article's lines are written, so no omission receipt of any kind accompanies this package. Citations: the retained excerpts carry no citation command at all, so no bibliography entry of the article is transcribed and no reference is redistributed. Reference adaptation: none was needed and none was made; every reference inside the delivered unit resolves inside it, so no unresolved reference or citation is produced. Printed slips kept literal: no printed word, symbol, hypothesis or number of the counted statement or of its proof was altered. Layout and class: the article's own class is \texttt{llncs} with packages including fontenc, graphicx, amssymb, amsmath, comment, xcolor, algpseudocode, algorithmicx, algorithm, ulem, rotating; the wrapper is the lane's pinned \texttt{amsart} editorial wrapper at 10pt on A4, which restates the theorem-like environments the retained text needs and adds the article's own macro definitions that the retained text needs (none), with extra wrapper packages (bm, bbm, dsfont, xspace, cancel, mathtools). The delivered numbering is the unit's own. Assets: no retained span contains a figure environment, a graphics-inclusion command, a PDF-page inclusion, a table or a TikZ picture; no image file is copied into this package, the package declares no asset dependency and no image file is redistributed with it. Licence: version arXiv:2510.01568v2 of this article is distributed under the Creative Commons Attribution 4.0 International licence, as recorded in the arXiv licence metadata for this exact version and shown on the abstract page https://arxiv.org/abs/2510.01568v2. A full rights notice is delivered at licenses/arxiv-2510.01568-CC-BY-4.0.txt. Characters: every retained excerpt is pure ASCII, so the delivered engine, XeLaTeX with its default Latin Modern text font, reports no missing character.
\begin{theorem}
If $p(x)$ is positive definite, and $\deg(p,x)=2d$, then  $p(x)$ can be expressed as the sum of squares of $d+1$ polynomials with decreasing degrees.  
If $p(x)$ is semi-positive definite, 
$\deg(p,x)=2d$,  
$$
p(x)=p_1(x)\times 
\prod_{j=1}^{k_0}(x-r_j)^{2}, \;
r_1,\ldots,r_{k_0}\in \mathbb{R},
$$
and $p_1(x)$ is positive definite, 
then $p(x)$ can be expressed as the sum of squares of $d-k_0+1$ polynomials with decreasing degrees from $d$ to $d-k_0$.   
\label{thm-2}
\end{theorem}

\begin{proof} First consider the case that $p(x)$ is positive definite. 
Without loss of generality, assume that the leading coefficient of $p(x)$ is 1. Since $p(x)$ is positive definite, $p(x)$ has $d$ pairs of complex conjugate roots  
\[
\alpha_i\pm\beta_i\sqrt{-1}, \quad i=1,2,\ldots,d,
\]
where $\alpha_i, \beta_i$ are real numbers, $\beta_i > 0$ for $1\leq i\leq d$.  
Thus, $p(x)$ can be rewritten as  
\begin{equation}
  p(x)=\prod_{i=1}^{d} \left[(x-\alpha_i)^2+\beta_i^2\right]=
\left(q_d(x)\vphantom{x_1^2}\right)^2+r_{d-1}(x),
\label{eq-0002}
\end{equation}
where $q_d(x)=\prod_{i=1}^{d}(x-\alpha_i)$ is a polynomial of degree $d$, and  
\begin{align}
r_{d-1}(x)&\;=\;
\sum_{k=1}^{d}\left(\frac{\beta_k}{x-\alpha_k}\,q_d(x)\right)^2\nonumber \\
&\;+\sum_{1\le k,l\le {d}, \ k\neq l}\left(\frac{\beta_k \beta_l}{(x-\alpha_k)(x-\alpha_l)}\,q_d(x)\right)^2
+\cdots+ \left(\prod_{k=0}^{d}\beta_k\right)^2.
\nonumber
%%%\label{eq-0003}
\end{align}
Clearly, $\deg(r_{d-1},x)\!=\!2(d-1)$, and the leading coefficient of $r_{d-1}(x)$ is 
$b_1 = \sum_{k=1}^d \beta_k^2 >0$. 
Since $r_{d-1}(x)$ is a sum of squares of polynomials, 
and $\beta_1,\ldots,\beta_d > 0$ ,  its constant term is positive, making $r_{d-1}(x)$ strictly positive definite. 

Applying the same argument we can write 
$r(x)$ in the following form:
\begin{equation}
r_{d-1}(x)=b_1\times \left(\left(q_{d-1}(x)\right)^2+r_{d-2}(x)\right)
\label{eq-rd1}
\end{equation}
where
$$
q_{d-1}(x)= \prod_{j-1}^{d-1}(x-\alpha'_j),\quad 
r_{d-2}(x)=
\sum_{j=1}^{d-1}\left(\frac{\beta'_k}
{x-\alpha'_k}q_{d-1}(x) \right)^2+\cdots+\prod_{j=1}^{d} {\beta'_j}^2,
$$ 
and $\alpha'_j,\beta'_j>0\,(j=1,\ldots,d-1)$ are real numbers, such that $\beta'_1,\ldots,\beta'_{d-1}>0$, and  
$$
\alpha'_j+\beta'_j\sqrt{-1}, \;
\alpha'_j-\beta'_j\sqrt{-1}, \quad j=1,2,\ldots,d-1
$$
are the $2(d-1)$ complex roots of $r_{d-1}(x)=0$. Apparently, $r_{d-2}$ is a positive definite of degree $d-2$ with leading coefficient $b_2=\sum_{j=1}^{d-1}{\beta'_j}^2>0$. Therefore, we can perform this computation for $r_{d-2}(x)$ and get
\begin{equation}
r_{d-2}(x)=b_2\times \left(
(\left(q_{d-2}(x)\right)^2+r_{d-3}(x)
\right),
\label{eq-rd2}
\end{equation}
where $q_{d-2}(x)$ is a polynomial of degree $d-2$ with leading coefficient $1$, and $r_{d-3}(x)$ is a positive definite polynomial of degree $2(d-3)$ with leading coefficient $b_3>0$. 
Inductively, we have
\begin{align}
r_{d-3}(x)=&\;b_3\times \left(
(\left(q_{d-3}(x)\right)^2+r_{d-4}(x)
\right),
\label{eq-rd3}\\
\vdots\;&\nonumber \\
r_{2}(x)=&\;b_{d-2}\times \left(
(\left(q_{2}(x)\right)^2+r_{1}(x)
\right),
\label{eq-rdj}
\end{align}
where $q_{j}\,(j=d-3,\ldots,2)$ are polynomial of degree $j$ with leading coefficient $1$, 
$r_{j}\,(j=d-4,\ldots,1)$ are  positive definite polynomial of degree $2j$ with leading coefficient $b_{d-j}>0$.  In particular,
\begin{equation}
r_1(x)=b_{d-1}\times (x^2+ax+b)
b_1(x)=b_{d-1}\times \left(\left(x+\frac{a}{2}\right)^2+
(b-\frac{a^2}{4})\right)
\label{eq-rdd}
\end{equation}
is a positive definite quadratic polynomial. Let $q_1(x)=x+a/2$, $b_d=b-a^2/4$. Then from \eqref{eq-0002} to \eqref{eq-rdd} we have
\begin{align}
p(x)=q_d(x)^2
+b_1\cdot q_{d-1}(x)^2
+b_1b_2\cdot q_{d-2}(x)^2\nonumber \\
+\cdots+b_1b_2\cdots b_{d-1}q_1(x)^2+
b_1b_2\cdots b_d,
\end{align}
which shows that $p(x)$ can be expressed as a sum of square of $d+1$ polynomials of degree from $d$ to zero, as claimed in the Theorem~\ref{thm-2}.  

Now, consider the case that 
$$
p(x)=p_1(x)\times r(x), \quad 
r(x)=(x-r_1)^2\cdots (x-r_{k_0})^2,
$$
where $p_1(x)$ is a positive definite polynomial of degree $2(d-k_0)$, $r_1,\ldots,r_{k_0}$ are real numbers.
Let
\[
p_1(x)=\left(\tilde{q}_{d-k_0}(x)\right)^2+\left(\tilde{q}_{d-k_0-1}(x)\right)^2+\ldots+\left(\tilde{q}_{0}(x)\right)^2,
\]
be the SOS of $p_1(x)$ into $d-k_0+1$ degree-descending-polynomials, 
where $\tilde{q}_{k}(x)$ (for $1\leq k\leq d-k_0$) is a real-coefficient polynomial of degree $k$ with a nonzero leading coefficient,  
and $\tilde{q}_{0}(x)$ is a nonzero real number. Let
$$
q_j=\tilde{q}_{j-k_0}\cdot r(x), \; \;
j=d,d-1,\ldots,k_0.
$$
Then
$$
p(x)=q_d(x)^2+q_{d-1}(x)^2+\ldots+
q_{k_0}(x)^2.
$$
which means that $p(x)$ can be written as the sum of squares of $d-k_0+1$ polynomials with decreasing degrees,  
where the highest and lowest degrees are $d$ and $k_0$, respectively.  

{Theorem~\ref{thm-2} is proved.}  
\qed  
    
\end{proof}  

\end{document}
